% `tex.web` §68's `print_ASCII` calls §59's `print`, and §59 has a case
% *before* any `^^` form:
%
%     if (@<Character |s| is the current new-line character@>) then
%       if selector<pseudo then
%         begin print_ln; return;
%         end;
%
% So the character equal to `\newlinechar` prints as an actual line break —
% in a box display, in a `short_display`, everywhere `print` is used.
%
% INITEX leaves `\newlinechar` at nought and `\defaulthyphenchar` at nought
% too, so a font loaded without one hyphenates with character 0 and TeX writes
% that hyphen as a line break.  That is what `\setbox0\hbox{A\-A}\showbox0`
% shows below: `..\r ` and then nothing, where we wrote `^^@`.
%
% It is not only that corner: LaTeX commonly sets `\newlinechar` to `^^J`, so
% any document that sets it and shows a box is affected.
%
% `detok_toks_string` already had the rule for `\message` and `\write`; this is
% the same rule reaching the printers that go through §59.
\catcode`\{=1 \catcode`\}=2 \catcode`\^=7
\font\r=cmr10 \r \hsize100pt \parindent0pt
\showboxbreadth99 \showboxdepth99 \tracingonline0
% \hyphenchar defaults to \defaulthyphenchar, which INITEX leaves at nought.
\setbox0\hbox{A\-A}\showbox0
% With \newlinechar elsewhere, character 0 prints in its `^^` form again.
\newlinechar=`\^^J
\setbox1\hbox{A\-A}\showbox1
% And the character that *is* \newlinechar prints as a break.
\catcode1=11
\setbox2\hbox{A^^AA}\newlinechar=1 \showbox2
\end
